% !TEX program = xelatex
\documentclass[11pt,a4paper,fontset=fandol]{ctexart}

\usepackage{amsmath,amssymb}
\usepackage{geometry}
\usepackage{booktabs}
\usepackage{array}
\usepackage{makecell}
\usepackage{float}
\usepackage{algorithm}
\usepackage{algpseudocode}
\usepackage{caption}
\usepackage{cite}
\usepackage[hidelinks]{hyperref}
\usepackage{titlesec}
\usepackage{enumitem}

\geometry{margin=2.2cm}
\setlength{\parindent}{0pt}
\setlength{\parskip}{0.25em plus 0.05em minus 0.03em}
\setlength{\textfloatsep}{6pt plus 2pt minus 2pt}
\setlength{\floatsep}{6pt plus 2pt minus 2pt}
\setlength{\intextsep}{6pt plus 2pt minus 2pt}
\setlength{\abovedisplayskip}{5pt plus 2pt minus 2pt}
\setlength{\belowdisplayskip}{5pt plus 2pt minus 2pt}
\titleformat{\section}{\large\bfseries}{{\thesection}}{0.5em}{}
\titleformat{\subsection}{\normalsize\bfseries}{{\thesubsection}}{0.5em}{}
\titleformat{\subsubsection}{\normalsize\bfseries}{{\thesubsubsection}}{0.5em}{}
\titlespacing*{\section}{0pt}{1.0ex plus 0.2ex minus 0.2ex}{0.5ex}
\titlespacing*{\subsection}{0pt}{0.8ex plus 0.2ex minus 0.2ex}{0.35ex}
\titlespacing*{\paragraph}{0pt}{0.4ex}{0.5em}
\setlist{nosep,leftmargin=1.5em}
\captionsetup{font=small,skip=4pt}
\renewcommand{\algorithmicrequire}{\textbf{输入:}}
\renewcommand{\algorithmicensure}{\textbf{输出:}}

\title{\bfseries 基于 CP 分解与数据感知逐层微调的 CNN 压缩方法研究}
\author{矩阵分析课程期末项目报告}
\date{\today}

\begin{document}

\maketitle

\begin{abstract}
深度卷积神经网络通常具有较高的参数量和推理计算量，这给资源受限场景下的部署带来困难。张量分解通过将权重张量重新参数化为低秩因子，为卷积层压缩提供了一种结构化方法。本文研究 CP 分解在 CNN 卷积层压缩中的应用：首先使用 CP-ALS 分解卷积核张量，并将原始卷积层替换为四个轻量卷积层；随后比较直接替换、全局微调、迭代式逐层微调和数据感知逐层微调四种压缩/恢复设置。实验在 MNIST 与 Fashion-MNIST 上进行，模型包括 LeNet-5 与 SmallCNN。结果显示，直接替换后的准确率接近随机；全局微调能够恢复部分性能，但在 Fashion-MNIST 上仍有明显精度损失；迭代式逐层微调显著改善恢复效果，而加入数据感知的层输出重构后，模型在保持 $51.44\%$（LeNet-5）和 $77.50\%$（SmallCNN）FLOPs 降低的同时，准确率基本恢复到稠密基线模型水平。

\textbf{关键词：} 模型压缩；CP 分解；CP-ALS；卷积神经网络；逐层微调
\end{abstract}

\section{引言}
近年来，卷积神经网络（CNN）在图像分类任务中取得了稳定而突出的表现，是深度学习方法的重要代表之一\cite{lecun2015deep}。但模型性能的提升通常伴随更大的参数规模和计算量；对于需要在边缘设备或低功耗环境中部署的模型而言，卷积层带来的存储与推理成本是一个重要限制。模型压缩的目标是在尽量保持预测性能的前提下，降低参数数量和计算复杂度。常见压缩方法包括剪枝、量化、知识蒸馏和低秩分解；其中，张量分解直接利用卷积核权重的多维结构，是一种具有明确数学形式的结构化压缩方法\cite{kolda2009tensor,comon2014tensors,panagakis2021tensor}。

CP 分解是基础的张量分解方法之一，它将一个高阶张量近似为若干秩一张量之和\cite{kolda2009tensor}。对于二维卷积层，其权重天然构成一个四阶张量，四个维度分别对应输出通道、输入通道、卷积核高度和卷积核宽度。已有研究将低秩分解用于卷积层加速，并进一步把完整四维卷积核的 CP 分解实现为若干小卷积层的串联结构\cite{denton2014exploiting,lebedev2015speeding}。若目标秩较小，该结构即可减少参数量与 FLOPs。

这一压缩过程的关键困难在于精度恢复。虽然 CP 分解能够减小权重重构误差，但该目标并不直接保证网络的中间特征和最终分类函数保持不变。直接替换卷积层通常会导致准确率大幅下降；一次性替换所有目标层后进行全局微调虽然简单，但多个层的近似误差会同时传播到后续网络，恢复效果有限。本文因此采用迭代式逐层微调，并进一步引入数据感知的层输出重构：每替换一层后，先使 CP 层在当前特征输入上的输出接近替换前的稠密层输出，再进行分类微调。实验在 MNIST 与 Fashion-MNIST、LeNet-5 与 SmallCNN 上验证该策略的有效性。

\section{模型与问题描述}
\subsection{卷积神经网络层}
CNN 中的一个标准二维卷积层由四阶权重张量
\[
\mathcal{W}\in\mathbb{R}^{O\times I\times H\times W}
\]
定义，其中 $O$ 和 $I$ 分别为输出和输入通道数，$H$ 和 $W$ 分别为卷积核高度和宽度。忽略偏置时，该层参数量为 $O I H W$。本文的目标是用参数更少的结构重新表示 $\mathcal{W}$，同时尽量保持原网络的分类性能。

\subsection{典范多元（CP）分解}
CP 分解将一个高阶张量近似表示为多个秩一张量之和。对卷积核张量而言，其形式为
\begin{equation}
\mathcal{W}[o,i,h,w]\approx\sum_{r=1}^{R}\lambda_r A[o,r]B[i,r]C[h,r]D[w,r],
\end{equation}
其中 $R$ 为 CP 秩，$A,B,C,D$ 分别对应输出通道、输入通道、卷积核高度和卷积核宽度四个模态。压缩后参数量约为 $R(O+I+H+W)$（不含偏置；精确值为 $R(O+I+H+W)+O$，偏置项保留在最后一个逐点卷积上）。

\subsection{问题描述}
按照该分解形式，原卷积层可替换为四个串联层：
\begin{enumerate}
  \item $1\times1$ 逐点卷积：$I\to R$；
  \item 深度可分离竖直卷积：$R\to R$，卷积核 $H\times1$；
  \item 深度可分离水平卷积：$R\to R$，卷积核 $1\times W$；
  \item $1\times1$ 逐点卷积：$R\to O$。
\end{enumerate}
该结构保持输入输出通道数和空间尺寸一致，同时将单个密集卷积拆解为计算成本更低的结构。

实验使用 MNIST 与 Fashion-MNIST，输入均为 $1\times28\times28$，类别数均为 10。本文考虑两个 CNN：
\begin{itemize}
  \item \textbf{LeNet-5}：两个卷积层（$1\to6$ 的 $5\times5$ 卷积，$6\to16$ 的 $5\times5$ 卷积）和三个全连接层，目标秩 $R=8$。
  \item \textbf{SmallCNN}：三个 $3\times3$ 卷积层（$1\to16\to32\to64$）和两个全连接层，目标秩 $R=16$。
\end{itemize}

\section{算法设计}
\subsection{基于 ALS 的 CP 分解}
CP 分解通过最小化 Frobenius 范数下的重构误差求解：
\begin{equation}
\min_{A,B,C,D}\left\|\mathcal{W}-\sum_{r=1}^{R}\lambda_r\,a_r\circ b_r\circ c_r\circ d_r\right\|_F^2 .
\end{equation}
该问题非凸。交替最小二乘法（ALS）每次固定除一个模态外的所有因子，将更新转化为线性最小二乘问题。设 $\mathbf{X}_{(n)}$ 为张量沿第 $n$ 个模态展开后的矩阵，$K^{(n)}$ 为其余因子的 Khatri--Rao 积，则更新为
\begin{equation}
U^{(n)} \leftarrow \mathbf{X}_{(n)}K^{(n)}\left((K^{(n)})^{\mathsf T}K^{(n)}\right)^\dagger .
\end{equation}
更新后对因子列归一化，并将尺度吸收到 $\lambda$ 中。

\begin{algorithm}[H]
\caption{CP-ALS 分解}
\begin{algorithmic}[1]
\Require 张量 $\mathcal{W}$，目标秩 $R$，迭代次数 $T$
\Ensure 权重 $\lambda$ 与因子矩阵 $A,B,C,D$
\State 随机初始化各因子，令 $\lambda\gets\mathbf{1}$
\For{$t=1$ to $T$}
  \For{每个模态 $n$}
    \State 计算其余因子的 Khatri--Rao 积 $K^{(n)}$
    \State 用最小二乘更新 $U^{(n)}$
    \State 归一化因子列并更新 $\lambda$
  \EndFor
\EndFor
\end{algorithmic}
\end{algorithm}

\subsection{微调策略}
CP 替换后模型性能会下降，因此需要分解后的适应过程。本文比较全局微调、迭代式逐层微调和数据感知逐层微调三种策略。

\subsubsection{全局微调}
全局微调是最直接的策略：先一次性分解并替换所有目标卷积层，再用较小学习率对整个模型做分类微调。其流程见算法 2。该方法实现简单，但网络需要同时适应多个层的结构变化。

\begin{algorithm}[H]
\caption{全局微调策略}
\begin{algorithmic}[1]
\Require 预训练模型 $M$，目标层列表 $L$，秩 $R$
\Ensure 压缩并微调后的模型 $M'$
\State $M'\gets M$
\For{每个目标层 $\ell\in L$}
  \State 使用 CP-ALS 分解 $\ell$ 的权重
  \State 用四层 CP 结构替换 $\ell$
\EndFor
\State 对 $M'$ 做全局分类微调
\end{algorithmic}
\end{algorithm}

\subsubsection{迭代式逐层微调}
实验发现，全局微调并不能完全恢复网络的性能，我们认为这可能是因为多个层同时被替换为低秩结构后，每层引入的近似误差会在前向传播中叠加，梯度信号需要同时补偿所有层的偏差，优化难度随层数增加而上升。
基于此，我们提出迭代式逐层微调方法：按深度逆序依次替换卷积层，优先训练更深层，每替换一层即立即进行分类微调，使网络在引入下一层低秩近似之前先适应当前的结构变化。
等到所有层都替换完成后，再进行一次全局微调以消除残余误差。

\subsubsection{数据感知逐层微调}
现有方法还存在一处优化目标的错位：CP-ALS 最小化的是权重的 Frobenius 重构误差，而非该层在真实样本上的输出误差。由于权重接近并不等价于输出接近，重构误差小的分解仍可能改变该层的输出行为。为此，数据感知逐层微调在每次替换后增加一个激活重构步骤，固定其余层参数，用真实训练数据驱动 CP 层，直接最小化替换前后该层在相同输入特征图上的输出差异。

设第 $\ell_k$ 层替换前的卷积映射为 $f^{\mathrm{ref}}_k$、替换后的 CP 映射为 $f^{\mathrm{CP}}_k$，激活重构目标为
\begin{equation}
\min_{\ell_k}\ \mathbb{E}_{x\sim\mathcal{D}}\left\|f^{\mathrm{CP}}_k(x_{\ell_k})-f^{\mathrm{ref}}_k(x_{\ell_k})\right\|_F^2,
\end{equation}
其中 $x_{\ell_k}$ 为当前模型（已部分替换）对输入 $x$ 前向传播至 $\ell_k$ 时的输入特征图。完成激活重构后再进行分类微调，使每层的替换在函数层面而非权重层面得到充分补偿。该策略的整体流程见算法 3。

\begin{algorithm}[H]
\caption{数据感知逐层微调策略}
\begin{algorithmic}[1]
\Require 预训练模型 $M$，目标卷积层列表 $L=(\ell_1,\dots,\ell_K)$（按深度升序排列），秩 $R$
\Ensure 压缩并微调后的模型 $M'$
\State $M'\gets M$（复制全部权重）
\For{$k=K,\,K{-}1,\,\ldots,\,1$（按深度逆序遍历）}
  \State $f^{\mathrm{ref}}_k\gets$ 当前 $M'$ 中层 $\ell_k$ 的卷积算子
  \State 对 $f^{\mathrm{ref}}_k$ 的权重做秩-$R$ CP-ALS 分解，将 $M'$ 中 $\ell_k$ 原地替换为四层 CP 结构
  \State \textbf{激活重构：}固定 $M'$ 中 $\ell_k$ 以外的所有参数，仅优化 $\ell_k$ 的 CP 权重，
  \Statex \quad\qquad 以训练集数据驱动，迭代 $T_{\mathrm{rec}}$ 个 epoch 最小化
         $\displaystyle\min_{\ell_k}\;\mathbb{E}_{x\sim\mathcal{D}}\!\left\|f^{\mathrm{CP}}_k(x_{\ell_k})-f^{\mathrm{ref}}_k(x_{\ell_k})\right\|_F^2$，
  \Statex \quad\qquad 其中 $x_{\ell_k}$ 为当前 $M'$ 对输入 $x$ 前向传播至 $\ell_k$ 时的输入特征图
  \State \textbf{分类微调：}以交叉熵损失对 $M'$ 全参数做 $T_{\mathrm{layer}}$ 个 epoch 的 Adam 微调
\EndFor
\State \textbf{最终全局微调：}以交叉熵损失对 $M'$ 全参数再做 $T_{\mathrm{final}}$ 个 epoch 的完整微调
\end{algorithmic}
\end{algorithm}

\section{数值结果}
\subsection{实验设置}
所有主实验使用 5 个随机种子，报告均值 $\pm$ 标准误差。基线模型训练 10 epoch，batch size 为 64，CP-ALS 迭代 50 次。LeNet-5 使用秩 $R=8$，SmallCNN 使用秩 $R=16$。逐层微调策略中，激活重构迭代次数 $T_{\mathrm{rec}}=1$，每层分类微调轮数 $T_{\mathrm{layer}}=2$，最终全局微调轮数 $T_{\mathrm{final}}=5$；全局微调策略亦采用 5 个 epoch。所有微调阶段均使用 Adam 优化器，逐层分类微调与激活重构学习率为 $10^{-3}$，全局微调学习率为 $10^{-4}$。

\subsection{性能指标}
本文使用三类指标评估压缩效果：测试集 Top-1 准确率、目标卷积层参数量，以及整网单次前向传播所需 FLOPs。参数统计包含偏置项；参数压缩比和 FLOPs 降低比例均相对于对应基线模型计算。

\begin{table}[H]
\centering
\caption{压缩统计}
\small
\begin{tabular}{lccccc}
\toprule
模型 & 秩 & 卷积参数 & 压缩比 & 基线模型 FLOPs & FLOPs 降低 \\
\midrule
LeNet-5 & 8 & $2572\to414$ & $6.21\times$ & 852.16 K & 51.44\% \\
SmallCNN & 16 & $23296\to2976$ & $7.83\times$ & 15619.97 K & 77.50\% \\
\bottomrule
\end{tabular}
\end{table}

\subsection{结果与分析}
\begin{table}[H]
\centering
\caption{不同恢复策略的准确率对比（\%，5 seed，均值 $\pm$ 标准误差）}
\small
\begin{tabular}{llccccc}
\toprule
数据集 & 模型 & 基线模型 & CP 无微调 & Global FT & Layer-wise FT & Data-aware FT \\
\midrule
MNIST & LeNet-5 & $98.89\pm0.06$ & $9.21\pm0.84$ & $96.95\pm0.24$ & $98.77\pm0.07$ & $\mathbf{99.02\pm0.02}$ \\
MNIST & SmallCNN & $99.08\pm0.06$ & $10.79\pm0.62$ & $97.96\pm0.14$ & $99.18\pm0.03$ & $\mathbf{99.33\pm0.02}$ \\
Fashion-MNIST & LeNet-5 & $90.17\pm0.14$ & $9.59\pm0.89$ & $83.89\pm0.71$ & $89.39\pm0.11$ & $\mathbf{90.31\pm0.08}$ \\
Fashion-MNIST & SmallCNN & $92.29\pm0.07$ & $9.49\pm0.28$ & $87.06\pm0.45$ & $92.03\pm0.04$ & $\mathbf{92.43\pm0.09}$ \\
\bottomrule
\end{tabular}
\end{table}

\paragraph{CP 无微调。}直接替换在所有设置下准确率均接近随机猜测（$\approx10\%$），标准误差较大（LeNet-5 可达 $\pm0.89\%$），说明 CP-ALS 所优化的权重重构误差与网络函数的保持之间存在根本性断层：低秩近似在 Frobenius 范数意义下可以做到较小，但叠加激活函数和多层传播后，输出分布的偏移足以使分类完全失效。

\paragraph{全局微调。}
全局微调在 MNIST 上恢复效果尚可（LeNet-5: $96.95\%$，SmallCNN: $97.96\%$），但在更难的 Fashion-MNIST 上出现明显的精度损失（LeNet-5 与基线差距达 $6.28\%$，SmallCNN 差距达 $5.23\%$），且标准误差显著偏大（LeNet-5 $\pm0.71\%$），表明优化过程不稳定。这印证了我们此前的假设:多层同时被替换后，梯度信号需同时补偿所有层的近似误差，优化路径更难收敛

\paragraph{逐层微调。}迭代式逐层微调在所有设置下均大幅优于全局微调。MNIST 上两个模型均接近或超过基线（LeNet-5: $98.77\%$，SmallCNN: $99.18\%$），Fashion-MNIST 上差距缩小至 $1\%$ 以内。标准误差相比全局微调明显下降，说明逐层分解带来了更稳定的优化轨迹。这验证了"将多层联合适应问题分解为逐层单任务"的有效性。

\paragraph{数据感知逐层微调。}在逐层微调基础上加入激活重构后，四组实验均达到或略超基线模型，其中 MNIST/SmallCNN 甚至高于基线（$99.33\%$ vs $99.08\%$），标准误差进一步压缩（最低达 $\pm0.02\%$）。精度提升的来源在于激活重构将 CP 分解的适应目标从权重空间的 Frobenius 误差转移到了真实输入特征图上的输出误差，使每层替换后的函数行为更接近原始稠密层，为后续分类微调提供了更好的初始点。上述四组实验同时保持表 1 中的参数压缩比和 FLOPs 降低。

\section{结论}
本文研究了基于 CP 分解的 CNN 卷积层压缩。实验表明，CP 分解可以显著降低卷积参数量与 FLOPs，但朴素替换会严重破坏模型精度；全局微调虽能恢复部分性能，但面对多层同时替换时恢复能力有限。逐层微调通过逐层引入结构变化，使模型逐步适应低秩卷积结构；进一步加入数据感知激活重构后，模型在 MNIST 与 Fashion-MNIST 上均基本恢复基线模型的准确率。

因此，基于 CP 分解的卷积层压缩不仅取决于张量近似本身，也取决于替换后的适应策略。对本实验规模而言，逐层微调和层输出重构是比单纯全局微调更稳定的精度恢复方案。后续可在更大规模网络和数据集上继续验证该结论，并研究每层秩的自动选择。

\section{分工与贡献}
本项目由两人合作完成，分工如下：
\begin{itemize}
  \item 周韧平：负责 CP-ALS 算法实现、全局微调和迭代式逐层微调的实验设计与实现， 贡献度：0.5。
  \item 陈驰：负责数据感知逐层微调的算法设计与实现，以及所有实验结果的统计与分析， 贡献度：0.5。
\end{itemize}
\begin{thebibliography}{99}

\bibitem{lecun2015deep}
Y. LeCun, Y. Bengio, and G. Hinton, ``Deep learning,'' \emph{Nature}, vol. 521, no. 7553, pp. 436--444, May 2015.

\bibitem{kolda2009tensor}
T. G. Kolda and B. W. Bader, ``Tensor decompositions and applications,'' \emph{SIAM Review}, vol. 51, no. 3, pp. 455--500, Aug. 2009.

\bibitem{comon2014tensors}
P. Comon, ``Tensors: A brief introduction,'' \emph{IEEE Signal Processing Magazine}, vol. 31, no. 3, pp. 44--53, May 2014.

\bibitem{panagakis2021tensor}
Y. Panagakis, J. Kossaifi, G. G. Chrysos, J. Oldfield, M. A. Nicolaou, A. Anandkumar, and S. Zafeiriou, ``Tensor methods in computer vision and deep learning,'' \emph{Proceedings of the IEEE}, vol. 109, no. 5, pp. 863--890, May 2021.

\bibitem{denton2014exploiting}
E. L. Denton, W. Zaremba, J. Bruna, Y. LeCun, and R. Fergus, ``Exploiting linear structure within convolutional networks for efficient evaluation,'' in \emph{Proc. Advances in Neural Information Processing Systems (NIPS)}, 2014.

\bibitem{lebedev2015speeding}
V. Lebedev, Y. Ganin, M. Rakhuba, I. Oseledets, and V. Lempitsky, ``Speeding-up convolutional neural networks using fine-tuned CP-decomposition,'' in \emph{Proc. International Conference on Learning Representations (ICLR)}, 2015.

\end{thebibliography}

\end{document}
